% BTS CIEL 1 — Réseaux — 07 — HTTP
% Compiler deux fois avec XeLaTeX depuis fichiers/.
\documentclass[aspectratio=169,11pt]{beamer}
\usepackage{fontspec}
\setmainfont{Latin Modern Sans}
\setsansfont{Latin Modern Sans}
\setmonofont{Latin Modern Mono}
\usepackage{polyglossia}
\setmainlanguage{french}
\usepackage{tikz}
\usetikzlibrary{arrows.meta,positioning,calc,shapes.geometric}
\usepackage{listings}
\usepackage{xcolor}
\definecolor{foret}{HTML}{1C322D}
\definecolor{creme}{HTML}{F8F3EE}
\definecolor{or}{HTML}{EBB552}
\definecolor{rose}{HTML}{F1CDBE}
\definecolor{vertclair}{HTML}{DDE9DF}
\definecolor{encre}{HTML}{24372E}
\setbeamercolor{normal text}{fg=creme,bg=foret}
\setbeamercolor{frametitle}{fg=or,bg=foret}
\setbeamercolor{title}{fg=creme}
\setbeamercolor{block title}{fg=foret,bg=or}
\setbeamercolor{block body}{fg=encre,bg=creme}
\setbeamercolor{alerted text}{fg=or}
\setbeamercolor{structure}{fg=or}
\setbeamertemplate{navigation symbols}{}
\setbeamertemplate{footline}{\hfill\color{or}\insertframenumber\hspace{0.5cm}\vspace{0.25cm}}
\setbeamertemplate{frametitle}{\vspace{0.35cm}\insertframetitle\par\vspace{0.15cm}\color{or}\rule{\linewidth}{0.6pt}}
\setlength{\parskip}{5pt}
\lstdefinestyle{webcode}{basicstyle=\ttfamily\small\color{encre},keywordstyle=\color{foret},stringstyle=\color{foret},
  backgroundcolor=\color{creme},frame=single,rulecolor=\color{or},framerule=1pt,
  xleftmargin=5pt,xrightmargin=5pt,columns=fullflexible,keepspaces=true,showstringspaces=false}
\newcommand{\question}[1]{\begin{beamercolorbox}[rounded=true,sep=12pt]{block body}\Large\textbf{#1}\end{beamercolorbox}}
\newcommand{\etiquette}[1]{{\color{or}\small\MakeUppercase{#1}}}
% Les schémas ont une surface claire, également lisible en noir et blanc.
\tikzset{
  boite/.style={draw=encre,fill=creme,text=encre,rounded corners,
    align=center,text width=2.8cm,minimum height=1cm},
  flux/.style={-{Stealth},thick,draw=encre},
  every node/.style={text=encre,font=\small}}
\newenvironment{schema}{%
  \begin{beamercolorbox}[sep=8pt,rounded=true]{block body}
  \centering\begin{tikzpicture}[>=Stealth]
}{\end{tikzpicture}\end{beamercolorbox}}
\title{Le protocole HTTP}
\subtitle{BTS CIEL 1 · Réseaux · 07 — HTTP}
\author{}
\date{}
\begin{document}

\begin{frame}[plain]
\vfill
\etiquette{Demander et recevoir une page Web}\par\bigskip
{\Huge\bfseries Le protocole HTTP}\par\bigskip
{\Large BTS CIEL 1 · Réseaux · 07 — HTTP}\par\bigskip
\begin{schema}
\node[boite,text width=3.4cm,minimum height=1.6cm,font=\large] (c) {Navigateur};
\node[boite,text width=3.4cm,minimum height=1.6cm,font=\large,right=3.5cm of c] (s) {Serveur Web};
\draw[flux] ([yshift=3mm]c.east)--node[above]{1. demande} ([yshift=3mm]s.west);
\draw[flux,dashed] ([yshift=-3mm]s.west)--node[below]{2. réponse} ([yshift=-3mm]c.east);
\end{schema}
\vfill
\end{frame}

\begin{frame}{Une page sur une autre machine}
\question{Comment demander la bonne page et connaître le résultat ?}
\begin{schema}
\node[boite,text width=3.4cm,minimum height=1.6cm,font=\normalsize] (a) at (0,2) {Client au comptoir};
\node[boite,text width=3.4cm,minimum height=1.6cm,font=\normalsize] (b) at (10,2) {Commerçant};
\draw[flux] ([yshift=3mm]a.east)--node[above=1mm]{1. une commande précise} ([yshift=3mm]b.west);
\draw[flux,dashed] ([yshift=-3mm]b.west)--node[below=1mm]{2. servi, rupture ou erreur} ([yshift=-3mm]a.east);
\node[font=\normalsize] at (5,0) {HTTP organise aussi demandes et réponses.};
\end{schema}

\end{frame}

\begin{frame}[shrink=20]{Lire une adresse Web}
\vspace{3mm}
\begin{schema}
\node[font=\normalsize] at (5,3.0) {Adresse Web = URL};
\node[font=\large,inner sep=0pt,anchor=base east] (proto) at (2.8,2.3) {\texttt{http}};
\node[font=\large,inner sep=0pt,anchor=base west] (sep) at (2.8,2.3) {\texttt{://}};
\node[font=\large,inner sep=0pt,anchor=base west] (host) at (sep.base east) {\texttt{www.example.com}};
\node[font=\large,inner sep=0pt,anchor=base west] (path) at (host.base east) {\texttt{/tarifs.html}};
\node[boite,text width=2.4cm,minimum height=1.2cm] (p) at (0,0) {protocole};
\node[boite,text width=4cm,minimum height=1.2cm] (h) at (5,0) {nom du serveur};
\node[boite,text width=3.2cm,minimum height=1.2cm] (c) at (10,0) {chemin demandé};
\draw (proto.south) -- (p.north);
\draw (host.south) -- (h.north);
\draw (path.south) -- (c.north);
\end{schema}
\begin{block}{Méthode}
1. Lisez le protocole avant \texttt{://}.\quad 2. Repérez le serveur.\\
3. Lisez le chemin à partir du premier \texttt{/} après le serveur.
\end{block}
\end{frame}

\begin{frame}{Une page arrive au navigateur}
\begin{schema}
\node[font=\normalsize] at (0,4) {Navigateur (client)};
\node[font=\normalsize,align=center] at (9,4) {Serveur\\\texttt{www.example.com}};
\draw[dashed] (0,3.5)--(0,.2); \draw[dashed] (9,3.5)--(9,.2);
\draw[flux] (0,3)--node[above,fill=creme,inner sep=2pt]{1. connexion TCP vers le port 80} (9,3);
\draw[flux] (0,2)--node[above,fill=creme,inner sep=2pt]{2. \texttt{GET /tarifs.html}} (9,2);
\draw[flux,dashed] (9,1)--node[above,fill=creme,inner sep=2pt]{3. \texttt{200 OK} + page} (0,1);
\end{schema}
Le navigateur affiche le contenu reçu. Exemple en HTTP/1.1.
\end{frame}

\begin{frame}{La requête précise la demande}
\begin{schema}
\node[boite,text width=3cm,minimum height=1.8cm,font=\normalsize] at (0,2) {\texttt{GET}\\méthode};
\node[boite,text width=3.5cm,minimum height=1.8cm,font=\normalsize] at (4.8,2) {\texttt{/tarifs.html}\\chemin};
\node[boite,text width=2.8cm,minimum height=1.8cm,font=\normalsize] at (9.5,2) {\texttt{HTTP/1.1}\\version};
\node[boite,text width=8cm,minimum height=1cm,font=\normalsize] at (4.8,0) {\texttt{Host: www.example.com}\\en-tête : nom du serveur};
\end{schema}
\begin{block}{Méthode}
1. Lisez la méthode et le chemin.\\
2. Repérez la version.\quad 3. Lisez l'hôte (nom du serveur) : \texttt{Host: www.example.com}.
\end{block}
\end{frame}

\begin{frame}{Consulter ou transmettre ?}
\begin{schema}
\node[boite,text width=4.5cm,minimum height=2.3cm,font=\large] at (0,0) {\texttt{GET}\\demander une page};
\node[boite,text width=4.5cm,minimum height=2.3cm,font=\large] at (6.5,0) {\texttt{POST}\\envoyer un formulaire};
\end{schema}
\begin{block}{Méthode}
1. Identifiez l'action.\quad 2. Consultez : pensez à \texttt{GET}.\\
3. Transmettez un formulaire : pensez à \texttt{POST}.
\end{block}
\end{frame}

\begin{frame}{La réponse indique le résultat}
\begin{schema}
\node[boite,text width=3.1cm,minimum height=2cm,font=\normalsize] at (0,2) {\texttt{HTTP/1.1}\\version};
\node[boite,text width=3.1cm,minimum height=2cm,font=\normalsize] at (4.5,2) {\texttt{200 OK}\\état : succès};
\node[boite,text width=3.1cm,minimum height=2cm,font=\normalsize] at (9,2) {Contenu\\page HTML};
\end{schema}
\begin{block}{Méthode}
1. Repérez la version.\quad 2. Lisez le code d'état.\\
3. Examinez le contenu reçu après les en-têtes.
\end{block}
\end{frame}

\begin{frame}{Lire les codes d'état}
\begin{beamercolorbox}[sep=12pt,rounded=true]{block body}
\centering\renewcommand{\arraystretch}{1.25}
\begin{tabular}{lll}
\textbf{Famille} & \textbf{Résultat} & \textbf{Exemple}\\\hline
2xx & Succès & 200 : réussi\\
3xx & Redirection & 301 : déplacement permanent\\
4xx & Erreur côté demande & 404 : ressource non trouvée\\
5xx & Erreur côté serveur & 500 : erreur interne\\
\end{tabular}
\end{beamercolorbox}
\begin{block}{Méthode}
1. Lisez le premier chiffre.\quad 2. Trouvez la famille.\\
3. Lisez le code complet pour préciser le résultat.
\end{block}
\end{frame}

\begin{frame}{HTTP en clair, HTTPS protégé}
Chiffrer : rendre le message illisible pour qui n'a pas la clé.
\begin{schema}
\node[boite,text width=2cm,minimum height=1cm] (a) at (0,3) {Navigateur};
\node[boite,text width=2cm,minimum height=1cm] (b) at (10,3) {Serveur};
\draw[flux] (a.east)--node[above=1mm]{HTTP, port 80} (b.west);
\draw (1.9,2.28) circle (0.1) (1.9,2.18) -- (1.9,1.92) (1.72,2.08) -- (2.08,2.08) (1.9,1.92) -- (1.75,1.78) (1.9,1.92) -- (2.05,1.78);
\draw[dotted,thick] (1.9,2.38)--(1.9,3);
\node[align=center] at (5.3,2.15) {Personne à l'écoute : \texttt{nom=Alice} lisible};
\node[boite,text width=2cm,minimum height=1cm] (c) at (0,0.8) {Navigateur};
\node[boite,text width=2cm,minimum height=1cm] (d) at (10,0.8) {Serveur};
\draw[flux] (c.east)--node[above=1mm]{HTTPS, port 443} (d.west);
\draw (1.9,0.08) circle (0.1) (1.9,-0.02) -- (1.9,-0.28) (1.72,-0.12) -- (2.08,-0.12) (1.9,-0.28) -- (1.75,-0.42) (1.9,-0.28) -- (2.05,-0.42);
\draw[dotted,thick] (1.9,0.18)--(1.9,0.8);
\node[align=center] at (5.3,-0.05) {Personne à l'écoute : \texttt{k8\#Qz\ldots}};
\end{schema}
TLS est le protocole qui protège les échanges HTTPS.
\end{frame}

\begin{frame}[shrink=20]{À vous : où chercher la panne ?}
\question{Pour chaque situation, où chercher la panne ?}
\begin{block}{Trois situations}
a. Le serveur répond \texttt{500}.\\
b. Le serveur répond \texttt{301} avec une nouvelle adresse.\\
c. Le serveur ne répond pas.
\end{block}
\begin{schema}
\foreach \y/\lab in {2.4/a,1.2/b,0/c} {
  \node[boite,text width=2.3cm,minimum height=0.7cm] (n\lab) at (0,\y) {Navigateur};
  \node[boite,text width=2.3cm,minimum height=0.7cm] (s\lab) at (10,\y) {Serveur};
  \draw[flux] ([yshift=2mm]n\lab.east)--([yshift=2mm]s\lab.west);
}
\draw[flux,dashed] ([yshift=-2mm]sa.west)--node[below]{a. \texttt{500}} ([yshift=-2mm]na.east);
\draw[flux,dashed] ([yshift=-2mm]sb.west)--node[below]{b. \texttt{301} + nouvelle adresse} ([yshift=-2mm]nb.east);
\node at (5,-0.25) {c. ?};
\end{schema}
\end{frame}

\begin{frame}[shrink=20]{Correction : interpréter les situations}
\vspace{3mm}
\begin{block}{Réponses}
a. \texttt{500} : cherchez le problème sur le serveur.\\
b. \texttt{301} : ce n'est pas une panne ; la page a déménagé.\\
c. Sans réponse, vérifiez le réseau (\texttt{ping}), puis que le serveur et son service HTTP sont lancés.
\end{block}
\begin{schema}
\foreach \y/\lab in {2.4/a,1.2/b,0/c} {
  \node[boite,text width=2.3cm,minimum height=0.7cm] (n\lab) at (0,\y) {Navigateur};
  \node[boite,text width=2.3cm,minimum height=0.7cm] (s\lab) at (10,\y) {Serveur};
  \draw[flux] ([yshift=2mm]n\lab.east)--([yshift=2mm]s\lab.west);
}
\draw[flux,dashed] ([yshift=-2mm]sa.west)--node[below]{a. \texttt{500} : serveur} ([yshift=-2mm]na.east);
\draw[flux,dashed] ([yshift=-2mm]sb.west)--node[below]{b. \texttt{301} : page déplacée} ([yshift=-2mm]nb.east);
\node at (5,-0.25) {c. réseau, serveur ou service arrêté};
\end{schema}
\end{frame}

\begin{frame}[shrink=20]{À vous : protéger un formulaire}
\question{Sur \texttt{https://www.example.org/compte}, quel port et quelle protection sont utilisés par défaut ?}
\begin{block}{À analyser}
Une personne à l'écoute du réseau peut-elle lire les données d'un formulaire transmis sur cette page ?
\end{block}
\begin{schema}
\node[boite,text width=2.4cm,minimum height=1.1cm] (n) at (0,2) {Navigateur};
\node[boite,text width=2.4cm,minimum height=1.1cm] (s) at (10,2) {Serveur};
\draw[flux] (n.east)--node[above]{\texttt{https://www.example.org/compte}} (s.west);
\draw[dotted,thick] (5,2)--(5,1.3);
\draw (5,1.1) circle (0.2) (5,0.9)--(5,0.3) (4.6,0.7)--(5.4,0.7) (5,0.3)--(4.7,0) (5,0.3)--(5.3,0);
\node[anchor=west] at (5.7,0.7) {Personne à l'écoute : ?};
\end{schema}
\end{frame}

\begin{frame}{Correction : une page HTTPS}
\begin{block}{Réponses}
Le port par défaut est \texttt{443}. HTTPS utilise TLS.\\
Les données transmises sont chiffrées : une personne à l'écoute ne peut pas les lire sans la clé.
\end{block}
\begin{schema}
\node[boite,text width=2.4cm,minimum height=1.1cm] (n) at (0,2) {Navigateur};
\node[boite,text width=2.4cm,minimum height=1.1cm] (s) at (10,2) {Serveur};
\draw[flux] (n.east)--node[above]{HTTPS, port 443, TLS} (s.west);
\draw[dotted,thick] (5,2)--(5,1.3);
\draw (5,1.1) circle (0.2) (5,0.9)--(5,0.3) (4.6,0.7)--(5.4,0.7) (5,0.3)--(4.7,0) (5,0.3)--(5.3,0);
\node[anchor=west] at (5.7,0.7) {Personne à l'écoute : \texttt{k8\#Qz\ldots}};
\end{schema}
\end{frame}

\begin{frame}{À retenir}
\begin{schema}
\node[boite,text width=3cm,minimum height=1.6cm,font=\normalsize] (a) {Navigateur};
\node[boite,text width=3cm,minimum height=1.6cm,font=\normalsize,right=4cm of a] (b) {Serveur};
\draw[flux] ([yshift=3mm]a.east)--node[above]{1. requête} ([yshift=3mm]b.west);
\draw[flux,dashed] ([yshift=-3mm]b.west)--node[below]{2. réponse + état} ([yshift=-3mm]a.east);
\end{schema}
\begin{block}{Quatre repères}
L'URL indique protocole, serveur et chemin.\\
HTTP/1.1 échange requêtes et réponses sur TCP, port 80 par défaut.\\
Le code d'état indique le résultat de la demande.\\
HTTPS utilise TLS ; son port par défaut est 443.
\end{block}
\etiquette{Chapitre suivant : DNS}
\end{frame}

\begin{frame}{Sources}
\begin{beamercolorbox}[sep=12pt,rounded=true]{block body}
\textbf{RFC 9110} — sémantique, méthodes et codes HTTP\\
\texttt{https://www.rfc-editor.org/rfc/rfc9110}\\[5pt]
\textbf{RFC 9112} — syntaxe et connexions HTTP/1.1\\
\texttt{https://www.rfc-editor.org/rfc/rfc9112}\\[5pt]
\textbf{RFC 8446} — TLS 1.3\\
\texttt{https://www.rfc-editor.org/rfc/rfc8446}\\[5pt]
Support enseignant : « TP DHCP, HTTP, DNS, FTP », p. 1--3 et 7--9.
\end{beamercolorbox}
\end{frame}
\end{document}
